% answers_short.tex —— 2013 年·上海卷·理科 · 简短答案速查
% 由 tools/make_short_answers.py 自 answers.tex 的 \gksolhead/\gkanswer 提取，
% 并与 2013/2013_数学真题_上海理_答案要点.md 逐题比对后生成，勿手改（改源 MD 后重跑）。

\begin{gknotice}
\textbf{使用说明}：本节只给最终答案便于快速核对，完整推导见解析版。答案为编写者独立推导，\textbf{非官方标准答案}。
\end{gknotice}

\gkpart{填空题}
\begin{itemize}[leftmargin=2.6em]
  \item[\textbf{1.}] $\dfrac13$。分子、分母同除以 $n$，极限为 $\dfrac{1}{3}$。
  \item[\textbf{2.}] $-2$。纯虚数要求 $m^2+m-2=0$ 且 $m^2-1\neq0$，故 $m=-2$。
  \item[\textbf{3.}] $0$。左式为 $x^2+y^2$，右式为 $-2xy$，故 $(x+y)^2=0$。
  \item[\textbf{4.}] $\arccos\left(-\dfrac13\right)$。由余弦定理比较得 $\cos C=-\dfrac13$。
  \item[\textbf{5.}] $-2$。通项中 $x^7$ 对应 $k=1$，系数 $5a=-10$。
  \item[\textbf{6.}] $\log_3 4$。令 $t=3^x$，方程化为 $t^2-2t-8=0$，取 $t=4$。
  \item[\textbf{7.}] $\dfrac{1+\sqrt5}{2}$。由 $\rho\cos\theta=1$ 与 $\rho=\cos\theta+1$ 得 $\rho^2-\rho-1=0$，取正根。
  \item[\textbf{8.}] $\dfrac{13}{18}$。总取法 $36$ 种，两球同为奇数的取法 $10$ 种，故所求为 $1-\dfrac{10}{36}$。
  \item[\textbf{9.}] $\dfrac{4\sqrt6}{3}$。取 $B(2,0)$，由 $BC=\sqrt2$、$\angle CBA=\dfrac\pi4$ 得 $C(1,1)$，代入得 $b^2=\dfrac43$，故 $2c=\dfrac{4\sqrt6}{3}$。
  \item[\textbf{10.}] $30d^2$。$19$ 个等距取值的方差为 $\dfrac{d^2(19^2-1)}{12}=30d^2$。
  \item[\textbf{11.}] $\dfrac23$。由 $\cos(x-y)=\dfrac12$ 及和差化积 $\sin2x+\sin2y=2\sin(x+y)\cos(x-y)$ 得 $\sin(x+y)=\dfrac23$。
  \item[\textbf{12.}] $\left(-\infty,-\dfrac87\right]$。$x>0$ 时 $f(x)=9x+\dfrac{a^2}{x}-7$，其最小值 $6|a|-7$ 需不小于 $a+1$，结合 $f(0)=0\geqslant a+1$ 得 $a\leqslant-\dfrac87$。
  \item[\textbf{13.}] $2\pi^2+16\pi$。截面面积分解为半径 $1$、长 $2\pi$ 的平放圆柱的截面 $4\pi\sqrt{1-y^2}$ 与底面积 $8\pi$、高 $2$ 的长方体截面 $8\pi$，由祖暅原理相加。
  \item[\textbf{14.}] $2$。由反函数对应关系，$[0,1)$ 映到 $(2,3]$，$[1,2)$ 映到 $[0,1)$，$[2,3]$ 映到 $[1,2]$；前两段无不动点，第三段只能在 $x=2$ 取等。
\end{itemize}

\gkpart{选择题}
\begin{center}\begin{tabular}{ccccc}
\toprule
题号 & 15 & 16 & 17 & 18 \\
\midrule
答案 & B & B & A & D \\
\bottomrule
\end{tabular}\end{center}

\gkpart{解答题}
\begin{itemize}[leftmargin=2.6em]
  \item[\textbf{19.}] $\dfrac23$。以 $A$ 为原点建系，平面 $D'AC$ 的法向量可取 $(1,-2,2)$，$B$ 到该平面的距离为 $\dfrac{|(1,-2,2)\cdot(2,0,0)|}{3}=\dfrac23$。
  \item[\textbf{20.}] (1) $[3,10]$。由 $200\left(5x+1-\dfrac3x\right)\geqslant3000$ 得 $x\geqslant3$，结合 $1\leqslant x\leqslant10$；(2) $x=6$ 千克/小时，最大利润 $457500$ 元。总利润 $90000\left(5+\dfrac1x-\dfrac{3}{x^2}\right)$，令 $u=\dfrac1x$，在 $u=\dfrac16$ 处取最大值
  \item[\textbf{21.}] (1) $\left(0,\dfrac34\right]$。需 $\left[-\dfrac{\pi}{4},\dfrac{2\pi}{3}\right]\subseteq\left[-\dfrac{\pi}{2\omega},\dfrac{\pi}{2\omega}\right]$，故 $\omega\leqslant\dfrac34$；(2) $\dfrac{43\pi}{3}$。零点为 $x=-\dfrac{\pi}{4}+k\pi$ 或 $x=\dfrac{5\pi}{12}+k\pi$，相邻零点间距交替为 $\dfrac{2\pi}{3}$ 与 $\dfrac{\pi}{3}$；含 $30$ 个零点的最短区间跨 $29$ 个间隔，取 $15$ 个 $\dfrac{\pi}{3}$ 与 $14$ 个 $\dfrac{2\pi}{3}$ 时最小
  \item[\textbf{22.}] (1) $x=-\sqrt3$（答案不唯一）。左焦点为 $(-\sqrt3,0)$，该竖直线与 $C_1$、$C_2$ 均有公共点；(2) 证明见解析。由 $|k||x|=|x|+1$ 得 $|k|>1$；过原点的直线若与 $C_2$ 有公共点则 $|k|>1$，此时代入 $C_1$ 无实根；$y$ 轴也与 $C_1$ 无公共点；(3) 证明见解析。圆内点 $P$ 满足 $|OP|<\dfrac1{\sqrt2}$；过 $P$ 的直线到原点距离小于 $\dfrac1{\sqrt2}$，再按 $|k|\leqslant1$ 与 $|k|>1$ 分别证明它不能同时与 $C_1$、$C_2$ 相交
  \item[\textbf{23.}] (1) $a_2=2$，$a_3=c+10$。$a_1+c=-2$ 落在中间段，$a_2=2$ 落在 $x\geqslant-c$ 段；(2) 证明见解析。分段计算得 $f(x)-x\geqslant c$ 对一切 $x$ 成立；(3) 存在，$a_1\geqslant-c$ 或 $a_1=-c-8$。等差数列公差 $d\geqslant c>0$，项最终进入 $x\geqslant-c$ 段，故 $d=c+8$；再逐项检验首项即可
\end{itemize}
